\documentclass{article}

\title{Szerzetesrendek}
\author{Lengyel Zoltán Péter, Simon László}
\date{Legutóbb frissítve: 2025. Március 10.}

\begin{document}

\maketitle

\section{Bencések}
\subsection{Pannonhalma}
\begin{itemize}
    \item 996. Szent Benedek követői alapították a monostort
    \begin{itemize}
        \item Szent Asztrik (megalapítja Pannonhalmát, neveli Szent Istvánt)
    \end{itemize}
    \item 15. században épült kerengő
    \begin{itemize}
        \item Szent Benedek Kápolna
    \end{itemize}
    \item Építészeti jellemzői
    \begin{itemize}
        \item Román stílus
        \begin{itemize}
            \item Lőrésszerű ablakok
            \item vastag falak (pl Jáki templom, Városligetben másolata)
            \item bélletes kapu (befelé szűkülő kapu)
            \item 1 főhajó + 1 mellékhajó
        \end{itemize}
        \item Korai gótika 
        \begin{itemize}
            \item Csúcsíves
            \item Támpillérek
            \item Rózsaablakok
            \item Környezeténél magasabb
            \item keskeny falak
            \item finom vonások
            \item pl Notre Dame
        \end{itemize}
        \item Késő gótika
        \begin{itemize}
            \item kevésbé magas
            \item vízköpő (díszelem a külső részeken)
        \end{itemize}
    \end{itemize}
    \item Apátság alapítólevele 1045.
\end{itemize}

\section{Ciszterciek}
\begin{itemize}
    \item Citeaux (Francia tartomány)> 
     \item 1098. Szent Róbert
     \item 1102. Szent Bernát - virágoztatja fel
     \item Kezdetben fehér $\rightarrow$ barna $\rightarrow$ fekete viselet (csuha)
     \item Carta Cartiatis (szeretet okirata) rend nűködésének regulája
     \item pl. Gödöllő
\end{itemize}

\section{Egyéb szerzetesrendek}
\begin{itemize}
    \item Premontreiek
    \item Karthauziak
    \item Ferencesek
    \item Domonkosok
    \item Johanniták
\end{itemize}

\end{document}